\section{A curvature obstruction to stable balanced-RM bundles in every
rank}\label{a-curvature-obstruction-to-stable-balanced-rm-bundles-in-every-rank}

Author: GPT-6 Astra Pro; inherited effort not exposed; HODGE\_DIVISION
\texttt{/root/sixfolds}. Date: 2026-11-21.

\textbf{Status:} SF-E811, SF-E811-RM and SF-E812 received whole-proof
mathematical PASS reviews from both \texttt{/root} and
\texttt{/root/imprimitive} on 2026-11-21 (GPT-6 Astra Pro; inherited
effort not exposed), after independent reconstruction of the complete
arguments. The root-requested rational-Kähler-class source restriction
is incorporated. The peer used the explicitly identified Siu existence
theorem and previously audited supplied S-21; no independent peer
reopening of Siu is claimed. The exact generic curvature certificate was
actually run and passed. This canonical proof concerns slope-stable or
slope-polystable \textbf{locally free} sheaves in the respective stated
scopes. It makes no assertion about arbitrary perfect complexes and does
not automatically extend to twisted sheaves.

\subsection{SF-E811. Strict mixed discriminant
positivity}\label{SF-E811}

Let \(X=V/\Lambda\) be a complex abelian fourfold with a
translation-invariant Kähler form

\[
 \omega=\omega_1+\omega_2,
\]

where \([\omega]\in H^2(X,\mathbf Q)\) is rational, \(V=V_1\oplus V_2\)
is an orthogonal complex splitting with \(\dim_{\mathbf C}V_i=2\), and
\(\omega_i\) is positive on \(V_i\) and zero on the other summand. Thus
a positive multiple of \([\omega]\) is a polarization, as required by
the exact projective existence source used below. Assume that neither of
the real four-form lines

\[
 \mathbf R\omega_1^2,\qquad \mathbf R\omega_2^2
 \quad\subset H^4(X,\mathbf R)
 \tag{E811.1}
\]

contains a nonzero rational class. For a holomorphic vector bundle \(E\)
of rank \(r>0\), set

\[
 \Delta(E)=c_2(E)-\frac{r-1}{2r}c_1(E)^2
           =\frac{c_1(E)^2}{2r}-\operatorname{ch}_2(E).
 \tag{E811.2}
\]

If \(E\) is slope-stable with respect to \(\omega\), then

\[
 \int_X\Delta(E)\,\omega_1\omega_2\ge0,
\]

and equality holds only when \(E\) is projectively flat. In particular
equality forces \(\Delta(E)=0\).

The proof first treats any Hermitian--Einstein connection. Its only use
of stability, apart from existence of this connection, is to make the
final finite deck action transitive on the parallel blocks. The parallel
algebra bundles used in the proof descend under the full holonomy; they
are not chosen summands that exist only on the universal cover.

\subsubsection{1. The mixed Chern--Weil
identity}\label{the-mixed-chernweil-identity}

The published Donaldson--Uhlenbeck--Yau existence theorem supplies a
Hermitian--Einstein Chern connection on \(E\). A complete projective
proof is in Siu, \emph{Lectures on Hermitian--Einstein metrics for
stable bundles and Kähler--Einstein metrics}, DMV Seminar 8 (1987),
chapter 1, §6.1--§6.4, author-reprint pp.35--39. We use this published
theorem as an input, not a new proof of the nonlinear existence theorem.
The normalization of a Hermitian--Einstein connection and its trace-free
part is also explicit in that source, chapter 1, §1.3--§1.4, author
pp.6--7.

Let \(F\) be the curvature, in the convention that its values on real
tangent vectors are skew-Hermitian, and put

\[
 F^0=F-\frac1r\operatorname{Tr}(F)I.
\]

Then \(F^0\) has type \((1,1)\), \(\Lambda_\omega F^0=0\), and the
Chern--Weil representative of (E811.2) is

\[
 \Delta(E)=\left[\frac1{8\pi^2}\operatorname{Tr}(F^0\wedge F^0)\right].
 \tag{E811.3}
\]

Here is a pointwise computation with all signs fixed. Use real unitary
coordinates \(x_j,y_j\), \(1\le j\le4\), and put

\[
 e_j=dx_j\wedge dy_j,\quad
 \omega_1=e_1+e_2,\quad\omega_2=e_3+e_4,\quad
 \mathrm{vol}=e_1e_2e_3e_4.
\]

A real \((1,1)\)-form can be written

\[
 \alpha=\sum_{j=1}^4 x'_j e_j+
 \sum_{i<j}\left(b_{ij}s_{ij}+c_{ij}t_{ij}\right),
\]

where

\[
 s_{ij}=dx_i\wedge dy_j+dx_j\wedge dy_i,
 \qquad t_{ij}=dx_i\wedge dx_j+dy_i\wedge dy_j.
\]

If \(\Lambda_\omega\alpha=0\), then \(\sum_jx'_j=0\), and direct
exterior multiplication gives

\[
 -\alpha^2\omega_1\omega_2
 =2\left((x'_1+x'_2)^2+
          \sum_{i\le2<j}(b_{ij}^2+c_{ij}^2)\right)\mathrm{vol}.
 \tag{E811.4}
\]

For example \(s_{ij}^2=t_{ij}^2=-2e_ie_j\); distinct off-diagonal basis
terms have zero cross pairing in this expression. The diagonal terms
give \(2(x'_1+x'_2)(x'_3+x'_4)\), and primitivity gives the displayed
sign. This proves (E811.4) directly; the exact script also verifies the
generic polynomial identity and its rank-nine positive semidefinite
matrix on the 15-dimensional primitive space.

Expand \(F^0\) in an orthonormal real basis of \(\mathfrak{su}(r)\) for
the inner product \(-\operatorname{Tr}(AB)\). Equation (E811.4), summed
over that basis, proves the nonnegativity in the theorem. If the
integral is zero, the pointwise nonnegative smooth density is zero
everywhere. Therefore

\[
 F^0=F_1+F_2,
 \qquad F_i\text{ has both real arguments in }V_i,
 \qquad\Lambda_{\omega_i}F_i=0.
 \tag{E811.5}
\]

In particular every mixed real curvature component vanishes. The
distinction between real and complex mixed components causes no gap:
type \((1,1)\) already removes the \((2,0)\) and \((0,2)\) terms, and
(E811.4) removes the remaining mixed terms.

\subsubsection{2. A product gauge and the labelled holonomy
algebras}\label{a-product-gauge-and-the-labelled-holonomy-algebras}

Pull back to the universal cover \(\widetilde X=V_1\times V_2\). The
smooth bundle is trivial there. In a unitary trivialization, subtract
the scalar part of the connection one-form. This produces a trace-free
unitary connection whose curvature is exactly \(F^0\). Choices of
determinant trivialization change the construction only by a scalar
one-form, which has no effect on conjugation on endomorphisms.

We prove the required splitting of this connection. Choose the frame
obtained by parallel transport first on the radial segment in \(V_1\)
and then on the radial segment in \(V_2\). Transport on each rectangle
made of one segment in each factor is independent of its order, because
the pulled-back connection on that rectangle has zero curvature by
(E811.5). Thus the frame is radial in both factors. If \(R_i\) is the
radial vector field of the \(i\)-th factor and \(A=A_1+A_2\) is the
connection one-form in this frame, then

\[
 A(R_1)=A(R_2)=0.
\]

For a constant vector \(u\in V_1\), the equation \(F^0(R_2,u)=0\)
becomes \(R_2(A_1(u))=0\). Smoothness at the origin of \(V_2\) implies
that \(A_1(u)\) is independent of the \(V_2\)-coordinate. The same
argument with the two factors exchanged shows that \(A_2\) depends only
on its own coordinate. The remaining mixed curvature equation then gives

\[
 A=A_1(z_1)+A_2(z_2),\qquad
 [A_1(z_1)(u_1),A_2(z_2)(u_2)]=0
 \tag{E811.6}
\]

for every choice of arguments. All these assertions concern real
coordinates; no holomorphic gauge is required.

Let \(H_i\) be the unitary holonomy subgroup generated by loops in the
\(i\)-th factor, based at a chosen point of \(\widetilde X\). Formula
(E811.6) implies that \(H_1\) and \(H_2\) commute and that the holonomy
of every loop on the universal cover lies in \(H_1H_2\). Let
\(\mathscr A_i\) be the unital complex \(*\)-algebra generated by
\(H_i\), as a subalgebra of the matrix algebra of the fiber. It is
finite-dimensional, and the two algebras commute.

These algebras define intrinsic \textbf{parallel subbundles of
\(\operatorname{End}(E)\) on \(X\)}. We spell out the descent. Changing
the base point along a path in the same factor conjugates its factor
holonomy group by parallel transport. Changing it along the other factor
identifies the group with that of the translated loops, by the
zero-curvature rectangle argument. A general path is handled by
(E811.6), and a loop on the universal cover normalizes both labelled
algebras because its holonomy belongs to the commuting product
\(H_1H_2\). Finally, every deck translation sends a loop tangent to
\(V_i\) to another loop with the same label \(i\); the bundle's deck
identification conjugates their parallel transports, up to the harmless
scalar from removing the determinant. Multiplying a generator by a
nonzero complex scalar does not change the unital complex algebra it
generates. Thus the construction is equivariant under every deck
transformation, and descends without choosing a lift of a projective
factor bundle or a root of the determinant.

The induced connections on \(\mathscr A_i\) are unitary for their
restricted Hilbert--Schmidt metrics. They are the Chern connections of
holomorphic vector bundles, since the induced \((0,2)\) curvature is
zero. Their curvatures are

\[
 F_{\mathscr A_i}=[F_i,-]|_{\mathscr A_i}.
 \tag{E811.7}
\]

Indeed the scalar part of \(F\) commutes with everything, and the other
factor's curvature commutes with \(\mathscr A_i\). The latter also
follows directly by differentiating the commuting factor parallel
transports. Equation (E811.5) makes these induced connections
Hermitian--Einstein with zero Einstein constant.

\subsubsection{3. Rationality kills the two adjoint
curvatures}\label{rationality-kills-the-two-adjoint-curvatures}

The curvature in (E811.7) has both arguments in the real four-plane
underlying \(V_i\). Its Chern--Weil degree-four form therefore belongs
pointwise to the one-dimensional space of real four-forms supported on
that plane. On a torus, averaging differential forms over translations
preserves their de Rham class. Averaging such a form leaves it in the
same constant one-dimensional space. Consequently

\[
 \operatorname{ch}_2(\mathscr A_i)\in
       \mathbf R\omega_i^2\cap H^4(X,\mathbf Q)=0
 \tag{E811.8}
\]

by the hypothesis (E811.1). Rationality is legitimate here:
\(\mathscr A_i\) is an actual finite-rank holomorphic vector bundle on
\(X\), not an unspecified projective or twisted bundle.

For a unitary primitive \((1,1)\) curvature \(G\), the Kähler Hodge-star
identity is

\[
 \operatorname{Tr}(G\wedge G)\wedge\omega^2
     =2\,|G|^2\,\frac{\omega^4}{4!}
 \tag{E811.9}
\]

with the positive inner product \(-\operatorname{Tr}\) and the usual
unitary norm convention. Equivalently, one can use only that the
constant on the right is positive; it follows by expanding in primitive
real \((1,1)\)-forms. Since
\(\operatorname{ch}_2=-\operatorname{Tr}(G^2)/(8\pi^2)\), (E811.8) and
(E811.9) force

\[
 F_{\mathscr A_1}=F_{\mathscr A_2}=0.
 \tag{E811.10}
\]

This conclusion kills the adjoint curvatures, not yet the centers. We
deal with the centers explicitly next.

\subsubsection{4. The central curvatures and a finite deck
orbit}\label{the-central-curvatures-and-a-finite-deck-orbit}

Every factor loop used to define \(H_i\) is contractible on \(X\), since
it is a closed loop on the simply connected universal cover. Its
holonomy on the flat bundle \(\mathscr A_i\) is therefore the identity.
On the matrix algebra this holonomy is conjugation by the corresponding
element of \(H_i\). Thus every element of \(H_i\) commutes with every
element of \(\mathscr A_i\). Since \(H_i\) generates \(\mathscr A_i\),
the algebra \(\mathscr A_i\) is commutative.

The unital \(*\)-algebra \(\mathscr A\) generated by \(\mathscr A_1\)
and \(\mathscr A_2\) is also commutative, because the two factor
algebras commute. It is a parallel bundle of finite-dimensional
commutative \(*\)-algebras. Such a fiber is a direct sum of copies of
\(\mathbf C\): simultaneous unitary diagonalization gives its primitive
orthogonal idempotents. Parallel transport permutes these finitely many
idempotents. A finite-index sublattice of \(\Lambda\) kills this
permutation action; for example, multiplication by \(r!\) suffices since
the number of nonzero idempotent blocks is at most \(r\).

Let \(p:X'\to X\) be the resulting finite étale isogeny. The pullback
splits into parallel holomorphic vector subbundles

\[
 p^*E=\bigoplus_{j=1}^s E_j.
 \tag{E811.11}
\]

The infinitesimal holonomy given by \(F_i\) belongs to \(\mathscr A_i\):
this follows by taking the infinitesimal holonomy of a small factor
rectangle. On each primitive block all elements of \(\mathscr A\) act as
scalars. Hence the full curvature of each \(E_j\) is scalar, including
the original scalar determinant curvature. Each \(E_j\) is therefore
projectively flat. This does not assert that its scalar curvature is
zero.

The finite deck group acts transitively on these blocks. If it had two
or more orbits, the direct sum of one orbit of blocks would descend to a
proper parallel holomorphic subbundle of \(E\). A parallel subbundle of
a Hermitian--Einstein bundle has the same slope as the whole bundle,
because the Einstein constant restricts to the same scalar. This would
contradict slope stability.

Deck transformations of an isogeny of complex tori are translations and
act trivially on rational cohomology. Transitivity therefore gives

\[
 \frac{c_1(E_j)}{\operatorname{rk}E_j}=\alpha
      \quad\text{for one class }\alpha\in H^2(X',\mathbf Q)
      \text{ and every }j.
 \tag{E811.12}
\]

The ranks are equal as well, though equality of normalized classes is
all that is needed. Since each block is projectively flat, Chern--Weil
gives

\[
 \operatorname{ch}(E_j)=\operatorname{rk}(E_j)e^{\alpha},
 \qquad \operatorname{ch}(p^*E)=r e^{\alpha}.
 \tag{E811.13}
\]

Thus \(p^*\Delta(E)=0\), and the injectivity of rational cohomological
pullback under a finite map gives \(\Delta(E)=0\).

One can also see directly that the entire original connection is
projectively flat. Write the scalar curvature of the block \(E_j\) as
\(\beta_j I\). Each \(\beta_j\) is closed of type \((1,1)\), and its
contraction with \(p^*\omega\) is the same constant. The Kähler
identities make it harmonic; this is Siu chapter 1, §1.4(ii), author
p.7. Equation (E811.12) makes these harmonic forms cohomologous, so they
are equal. Hence the curvature of \(p^*E\) is one scalar form times the
identity. The same property descends to \(E\). This handles all central
curvatures and finishes the proof of SF-E811.

\subsection{SF-E811-RM. Consequence for the remaining even-factor bundle
gate}\label{SF-E811-RM}

Suppose that \(X\) has genuine Rosati-invariant real multiplication by a
real quadratic field \(F\), with multiplicities \((2,2)\), and let
\(\theta=\theta_1+\theta_2\) be an \(F\)-compatible polarization. Then
no \(\theta\)-slope-stable vector bundle of positive rank \(N\) can have

\[
 e^B\operatorname{ch}(E)=N
     \left(1-\frac q2\theta_1^2\right)
     \left(1-\frac q2\theta_2^2\right),\qquad q>0,
 \tag{E811.14}
\]

for any rational algebraic divisor \(B\). The assertion holds for every
polarization degree, every rank and every integral order of real
multiplication. The presence of an additional finite translation/tensor
action does not create such an underlying stable bundle.

\textbf{Proof.} The positive invariant forms \(\theta_i\) give the
splitting in SF-E811. To check (E811.1), suppose that a nonzero rational
four-form were a real multiple of \(\theta_i^2\). That four-form is
decomposable, and its contraction kernel is exactly the other real
eigenspace. The kernel of a rational linear map is a rational subspace,
so that eigenspace would contain a nonzero rational vector. Choose
\(u\in F\setminus\mathbf Q\). The rational endomorphism induced by \(u\)
acts on this eigenspace by the irrational real scalar \(\sigma_j(u)\).
For a nonzero rational vector \(z\), the equation \(uz=\sigma_j(u)z\) is
impossible: any nonzero rational coordinate of \(z\) would make
\(\sigma_j(u)\) rational. This proves (E811.1), with no assumption that
\(X\) itself is simple.

The degree-two cohomology part of (E811.14) gives \(B=-c_1(E)/N\). Its
degree-four cohomology part then gives

\[
 \Delta(E)=\frac{Nq}{2}(\theta_1^2+\theta_2^2).
 \tag{E811.15}
\]

Because \(\theta_i^3=0\), this class pairs to zero with
\(\theta_1\theta_2\). By SF-E811 it must be zero. But its pairing with
\(\theta^2\) is

\[
 Nq\int_X\theta_1^2\theta_2^2
       =4Nq\int_X\frac{\theta^4}{24}>0,
\]

a contradiction. In fact only the normalized degree-four cohomology
component of (E811.14) is needed; the rest of the character is
irrelevant to this obstruction.

\subsection{SF-E812. Polystability cannot provide a weakly equivariant
semiregular
repair}\label{SF-E812}

In the balanced-RM setting of SF-E811-RM, let \(E\) be a
slope-polystable vector bundle with character (E811.14). Let
\(G\subset X\times\widehat X\) be any effective finite subgroup with a
specified coherent action and an equivariant structure on \(E\). Then
the full trace on invariant degree-two self-Ext is not injective. In
particular, no such polystable bundle is weakly \(G\)-semiregular,
regardless of its rank, endomorphism dimension, polarization degree or
group order.

This assertion permits actual polystable bundles with the stated Chern
character to exist. It excludes their proposed semiregularity property.
Its equivariant direct summands are constructed on \(X\) itself, not on
an auxiliary isogeny cover.

\subsubsection{1. All stable factors are projectively
flat}\label{all-stable-factors-are-projectively-flat}

Write the actual polystable decomposition

\[
 E=\bigoplus_j E_j,\qquad
 r_j=\operatorname{rk}E_j,\quad
 \alpha_j=c_1(E_j)/r_j,\quad
 \alpha=c_1(E)/N.
\]

The \(E_j\) are slope-stable vector bundles of the same
\(\theta\)-slope. Direct expansion of the discriminant gives

\[
 \Delta(E)=\sum_j\Delta(E_j)
       -\frac12\sum_j r_j(\alpha_j-\alpha)^2.
 \tag{E812.1}
\]

Represent \(\delta_j=\alpha_j-\alpha\) by its invariant real
\((1,1)\)-form. Equality of slopes says \(\Lambda_\theta\delta_j=0\).
Therefore each term

\[
 -\int_X\delta_j^2\theta_1\theta_2
\]

is nonnegative by the elementary identity (E811.4). Each
\(\int_X\Delta(E_j)\theta_1\theta_2\) is nonnegative by SF-E811. The
left side of (E812.1), paired with \(\theta_1\theta_2\), is zero by
(E811.15). Thus every nonnegative summand is zero. SF-E811 applied to
each stable factor shows that every \(E_j\) is projectively flat, and
hence

\[
 \operatorname{ch}(E_j)=r_j e^{\alpha_j}.
 \tag{E812.2}
\]

This step is compatible with the existence of polystable examples having
several different \(\alpha_j\): the negative variance term in (E812.1)
need not vanish as a cohomology class even though its mixed pairing
vanishes.

\subsubsection{2. The slope-class groups are actual equivariant
summands}\label{the-slope-class-groups-are-actual-equivariant-summands}

For each distinct rational class \(\beta\) among the \(\alpha_j\), let
\(E_\beta\) be the direct sum of the factors with \(\alpha_j=\beta\).
Polystable decomposition is unique at the level of stable isotypic
summands: nonisomorphic stable bundles of the same slope have no nonzero
Hom, and multiplicity spaces give the usual canonical isotypic sums.
Grouping these isotypic sums by \(\beta\) is therefore independent of a
choice of decomposition.

Translations and tensor products by degree-zero line bundles preserve
slope stability and the rational class \(c_1/r\). It follows that every
specified isomorphism

\[
 t_x^*E\otimes P\simeq E
\]

preserves each \(E_\beta\). The linearizations restrict, and their
coherence restricts as well. Consequently

\[
 E=\bigoplus_\beta E_\beta
 \tag{E812.3}
\]

is an actual direct-sum decomposition in the given equivariant category.
Each summand is nonzero and has

\[
 \operatorname{ch}(E_\beta)=r_\beta e^\beta,
 \qquad r_\beta>0.
 \tag{E812.4}
\]

\subsubsection{3. Trace/cap incompatibility for any one of these
summands}\label{tracecap-incompatibility-for-any-one-of-these-summands}

Suppose that the full trace on \(\operatorname{Ext}^2(E,E)^G\) were
injective. For the specified effective translation/tensor action,
supplied DEG\_L S-21 identifies the identity Hochschild sector with that
of \(X\), and its invariant Calabi--Yau adjunction makes degree-two
evaluation surjective. Thus the full trace image equals the cap-product
image of the whole character. This is precisely the implication used and
independently reviewed in SF-E802, including the equivariant setting.

Normalize all characters simultaneously by the same \(B\)-field shear in
(E811.14). In the Fock exterior model of SF-E801 the whole character
becomes, up to a nonzero scalar,

\[
 v=\omega'_1\wedge\omega'_2,
 \qquad W=U_1\oplus U_2,\quad\dim U_i=4,
\]

with each \(\omega'_i\in\bigwedge^2U_i\) nondegenerate. Its degree-two
cap kernel is

\[
 K=K_1\oplus K_2,
 \qquad K_i=\ker\left(\bigwedge^2U_i
                 \xrightarrow{\wedge\omega'_i}\bigwedge^4U_i\right),
 \qquad\dim K_i=5.
 \tag{E812.5}
\]

The common annihilator of \(K\) in \(\bigwedge^4W\) is exactly
\(\mathbf C v\), as proved in SF-E801.5. For clarity, the proof is
short: in bidegrees \((0,4),(1,3),(3,1),(4,0)\), multiplication by one
of the primitive spaces detects every nonzero component. In bidegree
\((2,2)\), annihilation by \(K_1\) restricts the first factor to
\(\mathbf C\omega'_1\), and annihilation by \(K_2\) restricts the second
to \(\mathbf C\omega'_2\). The nondegenerate exterior pairing on each
\(\bigwedge^2U_i\) proves these claims.

Because (E812.3) is an actual equivariant direct sum, the trace image of
each \(E_\beta\) is contained in the trace image of \(E\). Evaluation
followed by trace is cap product for each object. Hence

\[
 \operatorname{im}c^2_{e^B\operatorname{ch}(E_\beta)}
        \subset\operatorname{im}c^2_v.
 \tag{E812.6}
\]

The bilinear form \((z,w)\mapsto z\wedge w\wedge v\) on \(\bigwedge^2W\)
is symmetric, and the top exterior pairing is perfect. Taking
annihilators in (E812.6) gives

\[
 K\subset\ker c^2_{e^B\operatorname{ch}(E_\beta)}.
\]

Thus the normalized summand character belongs to the common-annihilator
line \(\mathbf C v\). But (E812.4) makes that character
\(r_\beta e^{B+\beta}\), a nonzero pure exponential. Its Fock
four-vector is decomposable, and its degree-two cap rank is six. The
balanced character has cap rank eighteen. They cannot be proportional.
This contradiction proves SF-E812.

No equivariant simplicity hypothesis was required. In particular, this
argument rules out weakly semiregular polystable repairs that would try
to exploit a large finite group order while leaving the actual
slope-class summands intact.

\subsection{Scope, sources and remaining
gates}\label{scope-sources-and-remaining-gates}

SF-E811 removes the \textbf{stable locally free} realization from the
remaining balanced-RM search even when the polarization/group factor is
even and the rank is at least four. SF-E812 removes every
\textbf{polystable locally free weakly equivariant semiregular}
realization. Neither assertion removes unstable locally free sheaves,
non-polystable semistable extensions, non-locally-free torsion-free
sheaves, arbitrary nonsimple indecomposable objects, twisted objects or
complexes with cohomology in several degrees. No simple perfect complex
is assumed to be slope-stable.

The only deep analytic input is the published Hermitian--Einstein
existence theorem. The exact accessible proof source is Yum-Tong Siu,
\emph{Lectures on Hermitian--Einstein metrics for stable bundles and
Kähler--Einstein metrics}, DMV Seminar 8, Birkhäuser, 1987, chapter 1,
§6.1--§6.4, \textbf{author-reprint pp.35--39}; the same author reprint
has §1.4(ii),(iii) on p.7 and Proposition 1.6 with proof on pp.8--9.
URL:
\url{https://people.math.harvard.edu/~siu/siu_reprints/dmv_book.pdf}.
The source clearly says that its proof presents Donaldson's improved
all-dimensional heat-flow argument. These are author-copy page numbers,
not an asserted identification with another printing.

The original published references are K. Uhlenbeck and S.-T. Yau,
\emph{On the existence of Hermitian--Yang--Mills connections in stable
vector bundles}, Comm. Pure Appl. Math. 39 (1986), S257--S293, DOI
\texttt{10.1002/cpa.3160390714}, with the 1989 note, 42 (1989),
703--707, DOI \texttt{10.1002/cpa.3160420505}; and S. K. Donaldson,
\emph{Infinite determinants, stable bundles and curvature}, Duke Math.
J. 54 (1987), 231--247, DOI \texttt{10.1215/S0012-7094-87-05414-7}.
Their publication metadata was checked. Direct PDFs of these originals
were unavailable in the present retrieval, so no original theorem number
or full original-proof read is claimed.

The mixed sum-of-squares identity, labelled algebra descent,
rational-volume-line argument, treatment of centers, and finite
deck-orbit argument are proved above. They invoke neither a preprint
theorem nor a conjectural geometric decomposition. The exact source
scope and the separate root/peer whole-proof checks are recorded in
SOURCE\_LOG.md and CHECK\_EVENTS.tsv. No formal proof-assistant
verification is claimed.
